\section{Experimental Design}
\label{sec:experiments}

We organize the experiments around four questions:
\begin{enumerate}
  \item Do the selection mechanisms produce distinct deletion sets under matched candidates and budgets?
  \item Do the selectors' reference or proxy objectives align with their stated targets?
  \item Relative to a budget-matched random request, how does a selected request change the measured GU outcome?
  \item How do these responses vary across GU methods, datasets, and attacker access settings, and how large is the signed gap from same-request retraining?
\end{enumerate}

\subsection{Ten-seed response-pattern case}
Figure~\ref{fig:response-pattern-profiles} uses the historical EXP-079/e79r5 batch: PubMed, split seed 2024, ten training seeds (42, 212, 2024, and 0--6), and deletion of 10\% of training-node candidates. Random, Degree, R-point, and D-full share the candidate pool and budget; the IF requests use a GCN selection source. Deletion removes training supervision and incident edges while retaining isolated feature rows; evaluation uses the original graph. Random's three request draws are averaged within each training seed; Degree reuses one deletion set across training seeds. Each GU endpoint is matched to full-graph GCN retraining on the same request. The partition-based methods are shard ensembles, so this reference does not isolate architecture-specific approximation error. The metric is raw test accuracy, without subtracting Random or Retrain. We show seed-level means and pointwise 95\% $t$ intervals, alongside each GU's pre-deletion mean where recorded; these exploratory intervals are not simultaneous tests of a ranking or equivalence. GIF and IDEA share a panel because their accuracy values coincide in this batch. GraphRevoker retains the archived GPA+optimal implementation. The batch is not pooled with the three-seed matrix, EXP-080, or the later revised sharded-method runs.

\subsection{Primary multi-method matrix}
The primary matrix covers node-deletion requests on Cora, CiteSeer, and PubMed with an OpenGU GCN backbone. We compare six graph-unlearning methods: GIF, GNNDelete, MEGU, IDEA, GraphEraser, and GraphRevoker \cite{wu2023gif,cheng2023gnndelete,li2024towards,dong2024idea,chen2022graph}. Each request removes 10\% of the eligible training-node candidates. The five selectors are Random, Degree, PageRank, R-point, and D-full. R-point is the practical influence-based setting used in the main-text comparison; D-full is a last-layer, two-hop setting included for reference.

The matrix contains 315 planned method conditions: 270 GU outputs and 45 full-retraining references. All 270 GU outputs have an identity-matched retraining result. We report three training seeds (42, 212, and 2024) as mean and sample standard deviation. In this matrix, Random, Degree, and PageRank reuse a selection across training seeds, whereas R-point and D-full selections vary with training seed. Thus, the reported spread includes model-training variability and, for the influence-based selectors, selection variability; it is not an estimate over independent random requests.

\draftnote{The registered split, candidate-pool construction, exact selector information access, training hyperparameters, and tie-breaking rules must be inserted from the final experiment manifests before submission.}

\subsection{Matched-reference analyses}
We add two current-parameter analyses to make the reference questions explicit. AAGU-032 contains 360 full-retraining conditions across ten selectors, four deletion budgets, three training seeds, and the three datasets. It compares selector outcomes under retraining; it does not measure GU behavior. AAGU-053 contains 48 strictly matched GNNDelete--retraining request pairs, with 16 requests per dataset at a 10\% budget and the training seed fixed to 42. The same selection identity is used at both endpoints.

These analyses use archived outputs from the 3000-epoch default configuration. The corresponding analysis report records that all AAGU-032 method outputs were cache hits; AAGU-053 includes both cache hits and producer executions. We therefore treat these as analyses of verified archived outputs, not as a claim that every condition was freshly recomputed in one run.

\subsection{Budget-response case study}
Figure~\ref{fig:exp080-budget-response} uses a separate historical batch, EXP-080/full-20261021. We fix CiteSeer, compare Random, Degree, R-point, D-full (\texttt{gt\_full}), and IM-RR, and vary deletion of training-node candidates over 5\%, 7.5\%, 10\%, 15\%, 22\%, 33\%, and 50\%. The six GU methods and explicit full-graph GCN retraining are paired on the actual deletion-node set within each budget, training seed, and request draw. The three training seeds are 42, 212, and 2024, with a fixed data-split seed of 2024. Random and Degree each reuse one deletion set across training seeds at a given budget; IM-RR uses three request draws (11, 22, 33), averaged within each training seed before computing seed means and sample SD. We report test accuracy and the percentage of fixed test nodes on which GU and same-request retraining predict different classes, with $\pm1$ sample SD bands rather than confidence intervals. The partition-based methods have different model families from the shared GCN retraining reference. GraphRevoker uses the GPA+optimal implementation in this archived batch; its aggregation implementation and training protocol predate the later revised sharded-method runs. We retain the original execution provenance and do not pool this case with those runs, the primary matrix, or the ten-seed illustration.

\subsection{Metrics and aggregation}
The primary matrix reports the F1 values recorded in its analysis report. We aggregate at the training-seed level and retain signed paired differences. For AAGU-032/053, the serialized utility field is named f1, but the output code computes test accuracy by comparing predicted and true labels on the test mask. We call that endpoint test accuracy throughout and keep it separate from the primary F1 tables.

For the AAGU-053 comparison, each dataset-level mean gives equal weight to its 16 matched requests. For the AAGU-032 selector comparison, we first average over three training seeds within each configuration and budget, then summarize the 12 configuration-budget groups per dataset. Those groups share configurations and are not independent samples. We use these results descriptively and do not report significance or equivalence from them.
